\documentclass[10pt,english,a4paper]{article}
\usepackage[utf8]{inputenc}
\usepackage[T1]{fontenc}
\usepackage{babel}

\usepackage[svgnames]{xcolor}
\usepackage{amsfonts}
\usepackage{a4wide}
\usepackage{graphicx}
\usepackage{url}
\usepackage{eurosym}
\usepackage{vmargin}
\usepackage{fancyhdr}
\usepackage{multicol}
\usepackage{lastpage}
\usepackage{eso-pic}
\usepackage{listings}
\usepackage{hyperref}
\usepackage{longtable}
\usepackage{syntax}
\usepackage{xspace}
\usepackage{booktabs}
\usepackage[capitalise,english,nameinlink]{cleveref} % load after algorithm2e and hyperref

\hypersetup{
  bookmarksdepth=2,
  bookmarksnumbered=true,
  bookmarksopen=true,
  bookmarksopenlevel=2,
  colorlinks=true,
  linktocpage=true,
  breaklinks=true,
  pageanchor=true,
  allcolors=[rgb]{0.3,0.3,0.99},
  pdftitle={Property Language Manual},
  pdfauthor={The MCC 2015-2019 formula team}
}

\setlength{\grammarindent}{11em}
\synctex=1

\newcommand\eqdef            {\mathrel{:=}}
\newcommand\tup[1]           {\langle#1\rangle}
\newcommand\set[1]           {{\{ #1 \mathclose \}}}
\newcommand\eg               {e.g.\@\xspace}
\newcommand\ie               {i.e.\@\xspace}
\newcommand\N                {\mathbb{N}}

\newcommand\ctla             {\textbf{A}\xspace}
\newcommand\ctle             {\textbf{E}\xspace}
\newcommand\ltlf             {\textbf{F}\xspace}
\newcommand\ltlg             {\textbf{G}\xspace}
\newcommand\ltlx             {\textbf{X}\xspace}
\newcommand\ltlu             {\textbf{U}\xspace}
\newcommand\logicnot         {\boldmath$\lnot$\xspace}
\newcommand\logicand         {\boldmath$\land$\xspace}
\newcommand\logicor          {\boldmath$\lor$\xspace}
\newcommand\quantexist       {\boldmath$\exists$\xspace}
\newcommand\quantall         {\boldmath$\forall$\xspace}
\newcommand\atomleq          {\boldmath$\leq$\xspace}
\newcommand\atomeq          {\boldmath$=$\xspace}
\newcommand\atomdead         {\textbf{deadlock}\xspace}
%\newcommand\atomql           {\textbf{quasilive}\xspace}
%\newcommand\atomstable       {\textbf{stablemarking}\xspace}
%\newcommand\atomlive         {\textbf{live}\xspace}
%\newcommand\atomsafe         {\textbf{safe}\xspace}
\newcommand\atomisfire[1]    {\textbf{is-fireable}(#1)}
\newcommand\atomplacebnd[1]  {\textbf{place-bound}(#1)}
\newcommand\atomtokenscnt[1] {\textbf{token-count}(#1)}
\newcommand\exprplus         {\boldmath$+$\xspace}
\newcommand\exprminus        {\boldmath$-$\xspace}

\newcommand{\CC}{%
  C\nolinebreak\hspace{-.05em}\raisebox{.4ex}{\scriptsize\bf +}\nolinebreak\hspace{-.10em}\raisebox{.4ex}{\scriptsize\bf +}%
}

%\definecolor{listinggray}{gray}{0.92}
\lstdefinelanguage{rnc}
{ morekeywords={element, complexType, sequence, ID, IDREF, complexContent, extension, attribute}
, morekeywords=[2]{name, minOccurs, maxOccurs, ref, type, substitutionGroup, base, abstract}
, morekeywords=[3]{xs, mcc}
, basicstyle=\scriptsize\fontfamily{fvm}\selectfont
, keywordstyle=\scriptsize\fontfamily{fvm}\selectfont\color{Purple}
, keywordstyle=[2]\scriptsize\fontfamily{fvm}\selectfont\color{DarkBlue}
, keywordstyle=[3]\scriptsize\fontfamily{fvm}\selectfont\color{ForestGreen}
, stringstyle=\scriptsize\fontfamily{fvm}\selectfont\color{FireBrick}
, morestring=[b]"
, backgroundcolor=\color{yellow!20}
}
\lstset{
  rulecolor=\color{black}
, breaklines=true
, tabsize=2
, flexiblecolumns=false
, keywordstyle=\bfseries
, basicstyle=\scriptsize\ttfamily
, xleftmargin=1cm
, xrightmargin=1cm
, language=rnc
}
\lstdefinelanguage{xsd}
{ morekeywords={element, complexType, sequence, ID, IDREF, complexContent, extension, attribute}
, morekeywords=[2]{name, minOccurs, maxOccurs, ref, type, substitutionGroup, base, abstract}
, morekeywords=[3]{xs, mcc}
, basicstyle=\scriptsize\fontfamily{fvm}\selectfont
, keywordstyle=\scriptsize\fontfamily{fvm}\selectfont\color{Purple}
, keywordstyle=[2]\scriptsize\fontfamily{fvm}\selectfont\color{DarkBlue}
, keywordstyle=[3]\scriptsize\fontfamily{fvm}\selectfont\color{ForestGreen}
, stringstyle=\scriptsize\fontfamily{fvm}\selectfont\color{FireBrick}
, morestring=[b]"
, backgroundcolor=\color{blue!10}
}
\lstdefinelanguage{mpl}
{ basicstyle=\scriptsize\fontfamily{fvm}\selectfont
, keywordstyle=\scriptsize\fontfamily{fvm}\selectfont\color{Purple}
, stringstyle=\scriptsize\fontfamily{fvm}\selectfont\color{FireBrick}
, morestring=[b]"
, backgroundcolor=\color{green!10}
}

\definecolor{rulercolor}{RGB}{255,255,255}
\definecolor{sectioncolor}{RGB}{0,0,0}
\definecolor{defaultcolor}{RGB}{11,71,50}

\oddsidemargin = 1cm
\evensidemargin = 1cm
\textwidth = 19cm
\topmargin = 1.2cm
\textheight = 24cm
\parskip = 0.1cm
\parindent = 0.0cm

%%%%%%%%%%%%%%%%%%%%%%%%%%%%%%%%%%%%%%%%%%%%%%%%%%%%%%
\fancyhf{}

\renewcommand{\headrulewidth}{0pt}
\renewcommand{\footrulewidth}{0pt}

\lhead{\textcolor{rulercolor}{Model Checking Contest @ Petri Nets}\\\ \\}
\chead{\textcolor{rulercolor}{Page {\thepage}/\pageref{LastPage}}\\\ \\}
\rhead{\textcolor{rulercolor}{Property Language Manual}\\\ \\}

\lfoot{}
\cfoot{}
\rfoot{}
%%%%%%%%%%%%% pour un fond
\newcommand\BackgroundPic{
\put(0,-160){
\parbox[b][\paperheight]{\paperwidth}{%
\vfill
\centering
\includegraphics[width=23.5cm,keepaspectratio]{backgroundNext.jpg}%
\vfill
}}}
%%%%%%%%%%%%%%%%%%%%%%%%% fin fond

\renewcommand{\rmdefault}{phv}
\renewcommand{\sfdefault}{cmr}
%\renewcommand{\ttdefault}{pcr}

%%%%%%%%%%%%%%%%%%%%%%%%%%%%%%%%%%%%
\newcommand{\fk}[1]{\color{blue}FK: #1\color{defaultcolor}\ }
\newcommand{\fp}[1]{\color{red}FK: #1\color{defaultcolor}\ }

\newcommand{\mcc}[0]{MCC}
\newcommand{\mccl}[0]{Model Checking Contest @ Petri Nets}
%%%%%%%%%%%%%%%%%%%%%%%%%%%%%%%%%%%%

\newcommand\myparagraph[1]{\color{sectioncolor}\smallskip\textbf{#1. }\color{defaultcolor}}
\newcommand\mysection[1]{\color{sectioncolor}\section{#1}\color{defaultcolor}}
\newcommand\mysubsection[1]{\color{sectioncolor}\subsection{#1}\color{defaultcolor}}

\pagestyle{fancy}
\AddToShipoutPicture{\BackgroundPic}

\begin{document}
\color{defaultcolor}

\vspace*{5cm}
\Huge
The Property Language Manual \\[10pt]
\Large
Model Checking Contest @ Petri Nets 2020

\vspace{3cm}
\large
Contact: \\[15pt]
The \mcc{} formula team, \\
\href{mailto:loig.jezequel@univ-nantes.fr}{Loïg Jezequel},
\href{mailto:emmanuel.paviot-adet@lip6.fr}{Emmanuel Paviot-Adet}
\vspace{2cm}

\vfill
~
\hfill
\scriptsize
v\@. Jan\@. 31, 2020
\normalsize

\clearpage
\tableofcontents

\clearpage
\mysection{Introduction}

\mysubsection{Scope}

This document is intended as a guide for the developers of
the tools participating in the \mccl{}, edition 2020.
It specifically contains:

\begin{itemize}
\item
  The definition of the property language used in the contest this year.
\item
  The definition of the language subset presented in each category of the
  contest.
%\item
%  Guidance on how to parse the XML files containing the properties.
\end{itemize}

\mysubsection{Categories of the contest}

The \mccl{} compares the competing verification tools on six
categories of verification goals:

\begin{enumerate}
\item State Space generation
\item Global properties analysis
\item Upper bound analysis
\item Reachability analysis
\item LTL analysis
\item CTL analysis
\end{enumerate}

To maximize tool participation, we have further divided some of these categories
into subcategories, where only restricted kinds of formulas
will be found:

\begin{center}
\begin{tabular}{lll}
\toprule
{\bf Category} & {\bf Subcategory}                       & {\bf Description} \\
\midrule
State Space generation & StateSpace & Counting the number of states in a marking graph \\
\midrule
Global properties analysis & GlobalProperties           & Existence of deadlock states \\
& QLDetection                 & Quasi-liveness detection \\
& StableMarkingDetection      & Detection of stable markings \\
& Liveness                    & Liveness detection \\
& 1Safe                       & 1-safeness detection \\
\midrule
Upper bound analysis & UpperBounds                    & Computing the exact upper bound of a set of places \\
\midrule
Reachability analysis & ReachabilityFireability        & Boolean combinations of checks of firability of transitions \\
& ReachabilityCardinality        & Boolean combinations of comparisons of numbers of tokens\\
\midrule
LTL analysis & LTLFireability                 & LTL with atomic propositions on firability of transitions \\
& LTLCardinality                 & LTL with atomic propositions on number of tokens\\
\midrule
CTL analysis & CTLFireability                 & CTL with atomic propositions on firability of transitions \\
& CTLCardinality                 & CTL with atomic propositions on numbers of tokens\\
\bottomrule
\end{tabular}
\end{center}

Each subcategory designates the
name of an XML file providing, for every benchmark, a set of formulæ for
tools competing in this subcategory to solve.
%
Most of the subcategories proposed this year were already proposed last
year. Only some global properties analysis are new.
%FIXME
%However, the class of formulas (\ie, the specific fragment of the
%BNF grammar presented in \cref{s:bnf}) contained in each of them has
%varied from last year.

Most of the formulas presented to tools have been randomly generated.
%Only very few of them were written by hand.
Unfortunately, most of the benchmarks submitted to the contest came without
example verification properties.
Writing them by hand is a very time consuming task (there are currently more
than 250 models!).
This is a very unfortunate situation but we have not found any way to
overcome it.

\iffalse
\mysubsection{Goals and evolution of the property language}

The property language for the \mccl{} is designed to maximize
the tools participating in the competition.
It is based on XML and allows to write structural, reachability, CTL, LTL
formulas.
It is tightly related to Petri nets, that are the modeling formalism of the
\mcc{}.

This language is designed to evolve in the future editions of the contest.
It is also designed to integrate with the \emph{Petri Net Markup Language}
(PNML) in the future.

The property language has been
extensively simplified with respect to the previous edition of the \mcc{}.
We kept compatibility in mind.
Any tool that participated in the previous edition of the contests should
only require minimal changes to accept the new version of the language.
\fi

\mysubsection{Outline}

This document presents the specification language in, conceptually, two steps.
\begin{itemize}
\item
  We first use a BNF grammar to describe the overall specification
  language.  This grammar is also used to formally define the class of
  formulas that will be presented to the tools participating on each
  subcategory of the contest (\cref{s:bnf}).
\item
  Second, the document defines the XML syntax of the files that will
  actually be passed to the tools (\cref{s:rng}).
\end{itemize}

%\Cref{s:parsing} provides some help on how to automatically generate code
%in various languages to parse the XML files containing the formulas.

\mysection{The BNF grammar of the property language}
\label{s:bnf}

Every property presented to a tool participating in any subcategory of the
contests can be derived from the following general grammar of the
property language:

\begin{grammar}
<formula> ::= <boolean-formula>
\alt <bound-formula>

<boolean-formula> ::= \ctla <boolean-formula>
\alt \ctle <boolean-formula>
\alt \ltlg <boolean-formula>
\alt \ltlf <boolean-formula>
\alt \ltlx <boolean-formula>
\alt <boolean-formula> \ltlu <boolean-formula>
\alt \logicnot <boolean-formula>
\alt <boolean-formula> \logicand <boolean-formula>
\alt <boolean-formula> \logicor <boolean-formula>
\alt \quantexist p <boolean-formula>
\alt \quantall p <boolean-formula>
\alt \quantexist t <boolean-formula>
\alt \quantall t <boolean-formula>
\alt \quantexist x <boolean-formula>
\alt \quantall x <boolean-formula>
\alt <atom>

<atom> ::= \atomdead
\alt \atomisfire{$t_1, \ldots, t_n$}
\alt <integer-expression> \atomleq <integer-expression>
\alt <integer-expression> \atomeq <integer-expression>

<integer-expression> ::=
     <integer-expression> \exprplus <integer-expression>
\alt <integer-expression> \exprminus <integer-expression>
\alt \textit{Integer constant}
\alt x
\alt \atomtokenscnt{$p_1, \ldots, p_n$}

<bound-formula> ::= \atomplacebnd{$p_1, \ldots, p_n$}

\end{grammar}

The property language considers two kinds of formulas:
\textit{Boolean formulas},
evaluating to a Boolean value, and \textit{bound formulas},
which evaluate to
%an interval $[a,b] \subseteq \N$ over the natural numbers.
a natural number.
Only the subcategory
\textit{UpperBounds} employs \textit{bound formulas}, all other
subcategories use only \textit{Boolean formulas}.

Operators \ctla and \ctle are the standard CTL operators.
Operators
\ltlg,
\ltlf,
\ltlx,
\ltlu
are the standard LTL operators.
We define that the \ltlx operator evaluate to false if no successor state
exists.  A \textit{Boolean formula} can also contain Boolean combinations of Boolean
formulas.
The \quantexist and \quantall operator quantify over places, transitions or an integer value $x$.
Atomic propositions are either
\begin{itemize}
\item
  \atomdead, asking that the current state enables no transition,
\item
  \atomisfire{$t_1, \ldots, t_n$}, which holds iff either $t_1$, or $t_2$,
  or \ldots or $t_n$ are enabled at the current state, or
\item
  (in)equalities of the form \synt{expr} \atomeq \synt{expr} or \synt{expr} \atomleq \synt{expr},
  where \synt{expr} is an integer expression.
\end{itemize}

Integer expressions can contain additions and subtractions of
\emph{Integer constants}, always taken from the set $\set{1, 2, 3}$, and
the following operator:
\begin{itemize}
\item
  \atomtokenscnt{$p_1, \ldots, p_n$}, which returns, for the current
  marking, the exact number of tokens contained in the set
  $\set{p_1, \ldots, p_n}$ of places.
\end{itemize}

\textit{Bound formulas} are formed only by atoms \atomplacebnd{$\cdot$}, whose semantics
is defined as follows:
\begin{itemize}
\item
  %Let $\N \eqdef \set{0, 1, 2, \ldots}$ be the set of natural numbers.
  The atom \atomplacebnd{$p_1, \ldots, p_n$} evaluates to some
%interval
%  $[a,b] \subseteq \N$ over the natural numbers,
  natural number $c$ so that $c$ is the maximum number of tokens that any reachable marking can put in all $n$
  places $p_1, p_2, \ldots, p_n$ at the same time.
%  where $a \in \N$ and $b \in \N \cup \set \infty$, if
%  the interval satisfies the following.
%  Let $c$ be the maximum number of tokens that any reachable marking can put in all $n$
%  places $p_1, p_2, \ldots, p_n$ at the same time.
%  Then~$c$ must be contained in the closed interval $[a,b]$.
\end{itemize}
%The interval $[a,b]$ is therefore an exact (if $a = b = c$) or imprecise
%(if $a \ne b$) overapproximation of the upper bound of the specified places.

Note that, in fact, in the current edition of the \mcc{} no subcategory of the
contest will contain formulas using additions and subtractions of integer
expressions. These operators remain here only for future use.

% FIXME
% give formal semantics to these formulae

This grammar defines a large class of formulæ.
Each subcategory of the contest contains formulas of only some specific
fragment of this grammar.
The following subsections define the specific syntax of the formul\ae{}
that tools participating in a given subcategory shall expect to be provided
with.

\Cref{s:rng} presents the correspondence between the above BNF grammar
and the syntax of the XML files that will actually be passed to the tools.

\mysubsection{Subcategory \textit{GlobalProperties}}

Tools participating in this subcategory will only find one formula to solve
for every benchmark:

\begin{grammar}

<formula> ::= <boolean-formula>

<boolean-formula> ::= \ctle \ltlf \atomdead
%\alt A G $\lnot$ deadlock

\end{grammar}

\mysubsection{Subcategory \textit{QLDetection}}

Tools participating in this subcategory will only find one formula to solve
for every benchmark:

\begin{grammar}

<formula> ::= <boolean-formula>

<boolean-formula> ::= \quantall t~ \ctle \ltlf \atomisfire{t}
%\alt A G $\lnot$ deadlock

\end{grammar}

\mysubsection{Subcategory \textit{StableMarkingDetection}}

Tools participating in this subcategory will only find one formula to solve
for every benchmark:

\begin{grammar}

<formula> ::= <boolean-formula>

<boolean-formula> ::= \quantexist p~ \quantexist x~ \ctla \ltlg \atomtokenscnt{p} \atomeq x
%\alt A G $\lnot$ deadlock

\end{grammar}

\mysubsection{Subcategory \textit{Liveness}}

Tools participating in this subcategory will only find one formula to solve
for every benchmark:

\begin{grammar}

<formula> ::= <boolean-formula>

<boolean-formula> ::= \quantall t~ \ctla \ltlg \ctle \ltlf \atomisfire{t}
%\alt A G $\lnot$ deadlock

\end{grammar}

\mysubsection{Subcategory \textit{1Safe}}

Tools participating in this subcategory will only find one formula to solve
for every benchmark:

\begin{grammar}

<formula> ::= <boolean-formula>

<boolean-formula> ::= \quantall p~ \ctla \ltlg \atomtokenscnt{p} \atomleq 1
%\alt A G $\lnot$ deadlock

\end{grammar}

\mysubsection{Subcategory \textit{UpperBounds}}

Observe that this is the only subcategory where formulas evaluate to something
else than a Boolean value, in particular to %an interval $[a, b] \subseteq \N$
%over the
a natural number. See the beginning of this section for the semantics
of the atom \atomplacebnd{$\cdot$}.

\begin{grammar}
<formula> ::= <bound-formula>

<bound-formula> ::= \atomplacebnd{$p_1, \ldots, p_n$}
\end{grammar}

\mysubsection{Subcategories \textit{ReachabilityFireability} and \textit{ReachabilityCardinality}}

Both subcategories contain formulas produced using the same grammar, the
only difference being the atomic propositions allowed in each one.
The grammar is the following:

\begin{grammar}
<formula> ::= <boolean-formula>

<boolean-formula> ::= \ctle \ltlf <state-formula>
\alt \ctla \ltlg <state-formula>

<state-formula> ::=  \logicnot  <state-formula>
\alt <state-formula> \logicand <state-formula>
\alt <state-formula> \logicor  <state-formula>
\alt <atom>
\end{grammar}

Now, the subcateory \emph{ReachabilityFireability} only contains
\atomisfire{$\cdot$} atoms:

\begin{grammar}
<atom> ::= \atomisfire{$t_1, \ldots, t_n$}
\end{grammar}

while the formulas in the \emph{ReachabilityCardinality} subcategory
uniquely contain integer inequalities:

\begin{grammar}
<atom> ::= <integer-expression> \atomleq <integer-expression>

<integer-expression> ::= \textit{Integer constant}
\alt \atomtokenscnt{$p_1, \ldots, p_n$}
\end{grammar}

Recall that the \textit{Integer constant} will be a number in $\set{1,2,3}$.


\mysubsection{Subcategories \textit{LTLFireability} and \textit{LTLCardinality}}

Both subcategories contain formulas produced using the same grammar, the
only difference being the atomic propositions allowed in each one.
The grammar is the following:

\begin{grammar}

<formula> ::= <boolean-formula>

<boolean-formula> ::= \ctla <path-formula>

<path-formula> ::=
     \ltlg <path-formula>
\alt \ltlf <path-formula>
\alt \ltlx <path-formula>
\alt <path-formula> \ltlu <path-formula>
\alt \logicnot <path-formula>
\alt <path-formula> \logicand <path-formula>
\alt <path-formula> \logicor <path-formula>
\alt <atom>

\end{grammar}

Now, the subcategory \emph{LTLFireability} only contains
\atomisfire{$\cdot$} atoms:

\begin{grammar}
<atom> ::= \atomisfire{$t_1, \ldots, t_n$}
\end{grammar}

while the formulas in the \emph{LTLCardinality} subcategory only contain
integer inequalities:

\begin{grammar}
<atom> ::= <integer-expression> \atomleq <integer-expression>

<integer-expression> ::= \textit{Integer constant}
\alt \atomtokenscnt{$p_1, \ldots, p_n$}
\end{grammar}


\mysubsection{Subcategories \emph{CTLFireability} and \emph{CTLCardinality}}

Both subcategories contain formulas produced using the same grammar, the
only difference being the atomic propositions allowed in each one.
The grammar is the following:

\begin{grammar}
<formula> ::= <boolean-formula>

<boolean-formula> ::=
     \ctla <path-formula>
\alt \ctle <path-formula>
\alt \logicnot <boolean-formula>
\alt <boolean-formula> \logicand <boolean-formula>
\alt <boolean-formula> \logicor <boolean-formula>
\alt <atom>

<path-formula> ::=
     \ltlg <boolean-formula>
\alt \ltlf <boolean-formula>
\alt \ltlx <boolean-formula>
\alt <boolean-formula> \ltlu <boolean-formula>
\end{grammar}

The subcategory \emph{CTLFireability} only contains \atomisfire{$\cdot$}
atoms:

\begin{grammar}
<atom> ::= \atomisfire{$t_1, \ldots, t_n$}
\end{grammar}

while the formulas in the \emph{CTLCardinality} subcategory contain
integer inequalities:

\begin{grammar}
<atom> ::= <integer-expression> \atomleq <integer-expression>

<integer-expression> ::= \textit{Integer constant}
\alt \atomtokenscnt{$p_1, \ldots, p_n$}
\end{grammar}

% NOTE - doesn't make sense to say this, we need to filter out in all
% categories in fact!!

% Observe that under the previous definitions it is possible to derivate a
% formula without any CTL operator, by using the last grammar production for
% the non-terminal \synt{boolean-formula}.
% However, such formul{\ae} will never be presented to a tool participating
% in these subcategories, we will ensure that
% \emph{formul{\ae} contain at least one CTL (either \ctla or \ctle)
% operator}.

\mysection{XML format of the property language}
\label{s:rng}

During the competition, tools will be provided, for every benchmark and
every subcategory of the contest, with one XML file containing the set of
formulas to solve for that subcategory and that benchmark.

The purpose of this section is defining the XML Schema describing the
structure of the XML files containing the formulas.
In fact, we describe a RelaxNG representation of the XML Schema.

Each XML file will be provided together with a plain text file describing
the formula in a more readable way.
The format of this text file is not guaranteed to be preserved in future
editions of the contest. This file is just provided for readability.

As an introducing example, consider the following property, given in XML
(left) and the corresponding to a set of formulas containing only one
formula (right) generated by the grammar in \cref{s:bnf}.

\begin{multicols}{2}
\begin{lstlisting}
   <?xml version="1.0"?>
   <property-set xmlns="http://mcc.lip6.fr/">
    <property>
     <id>Dekker-PT-010-test1</id>
     <description>
      Automatically generated formula.
     </description>
     <formula>
      <exists-path>
       <until>
        <before>
         <negation>
          <integer-le>
           <integer-constant>1</integer-constant>
           <tokens-count>
            <place>p3_0</place>
           </tokens-count>
          </integer-le>
         </negation>
        </before>
        <reach>
         <is-fireable>
          <transition>withdraw_4_8</transition>
         </is-fireable>
        </reach>
       </until>
      </exists-path>
     </formula>
    </property>
   </property-set>
\end{lstlisting}
\columnbreak
\vspace*{2cm}
\ctle ((\logicnot (1 \atomleq \atomtokenscnt{p3\_0})) U \atomisfire{withdraw\_4\_8})
\end{multicols}

What follows is the description of the RelaxNG grammar describing the XML
files provided to the tools.

\mysubsection{Property sets}
The \lstinline[language=xsd]!property-set! element is the root of the XML representation.
It contains one or more properties.
\begin{lstlisting}
  default namespace = "http://mcc.lip6.fr/"
  start = property-set

  property-set = element property-set {
    property*
  }
\end{lstlisting}

\mysubsection{Properties}
A property is composed of three mandatory parts: a unique identifier, a textual description
of the property, and the formula itself.
%It has a set of tags to give a hint about the class of tools that can compute the formula
%(structural, reachability, CTL or LTL). These hints do \emph{not} say that the formula
%is expressed exactly in the language given by the hint, but tells that given minor changes,
%it could be.

%A property also contains an optional part for the expected result.
%It gives the expected value and a textual explanation.
%% \begin{lstlisting}
%%   property = element property {
%%     element id {
%%       xsd:ID
%%     } &
%%     element description {
%%       text
%%     } &
%%     element tags {
%%       element is-structural { xsd:boolean } &
%%       element is-reachability  { xsd:boolean } &
%%       element is-ctl { xsd:boolean } &
%%       element is-ltl { xsd:boolean }
%%     } &
%%     element expected-result {
%%       element value {
%%         xsd:integer | xsd:boolean
%%       } &
%%       element explanation {
%%         text
%%       }
%%     }? &
%%     element formula {
%%       formula
%%     }
%%   }
%% \end{lstlisting}

\begin{lstlisting}
  property = element property {
    element id {
      xsd:ID
    } &
    element description {
      text
    } &
    element formula {
      formula
    }
  }
\end{lstlisting}


\mysubsection{Formul{\ae}}
Formul{\ae} are the body of properties.
They define what is expected to hold on the model.
Formul\ae{} are currently of two main types: formul{\ae} that return integers, and formul{\ae}
that return Booleans.

In the following, for each rule of the grammar defining formulæ we give the
RelaxNG representation (on the left) as well as the corresponding part of
the BNF grammar (on the right).

\begin{multicols}{2}
\begin{lstlisting}
  formula =
      boolean-formula
    | bound-formula
\end{lstlisting}
\columnbreak
\begin{grammar}
<formula> ::= <boolean-formula>
\alt <bound-formula>
\end{grammar}
\end{multicols}

\subsubsection{CTL state operators}

These operators correspond to the \ctla and \ctle operators of CTL.

\begin{multicols}{2}
\begin{lstlisting}
  boolean-formula =
      ...
    | element all-paths {
        boolean-formula
      }
    | element exists-path {
        boolean-formula
      }
    | ...
\end{lstlisting}
\columnbreak
\begin{grammar}
<boolean-formula> ::= \ctla <boolean-formula>
\alt \ctle <boolean-formula>
\alt \ldots
\end{grammar}
\end{multicols}

\subsubsection{Path operators}

These operators are the
\ltlg,
\ltlf,
\ltlx, and
\ltlu
operators of CTL and LTL
Recall that the \lstinline!next! operator should evaluate to false if no
successor state exists.

\begin{multicols}{2}
\begin{lstlisting}
  boolean-formula =
      ...
    | element globally {
        boolean-formula
      }
    | element finally {
        boolean-formula
      }
    | element next {
        boolean-formula &
    | element until {
        element before {
          boolean-formula
        } &
        element reach {
          boolean-formula
        } &
      }
    | ...
\end{lstlisting}
\columnbreak
\setlength{\grammarindent}{10em}
\begin{grammar}
<boolean-formula> ::=
     \ltlg <boolean-formula>
\alt \ltlf <boolean-formula>
\alt \ltlx <boolean-formula>
\alt <boolean-formula> \ltlu <boolean-formula>
\alt \ldots
\end{grammar}
\end{multicols}

\subsubsection{Boolean operators}

These are usual Boolean operators.

\begin{multicols}{2}
\begin{lstlisting}
  boolean-formula =
      ...
    | element negation {
        boolean-formula
      }
    | element conjunction {
        boolean-formula,
        boolean-formula+
      }
    | element disjunction {
        boolean-formula,
        boolean-formula+
      }
    | ...
\end{lstlisting}
\columnbreak
\setlength{\grammarindent}{6em}
\begin{grammar}
<boolean-formula> ::=
     \logicnot <boolean-formula>
\alt <boolean-formula> \logicand <boolean-formula>
\alt <boolean-formula> \logicor <boolean-formula>
\alt \ldots
\end{grammar}
\end{multicols}


\subsubsection{Atomic propositions}

There is three types of atomic propositions.
Recall that
\lstinline!deadlock! evaluates to true if the current state is a deadlock
(has no successor) and
\lstinline!is-fireable! evaluates to true if one of the set of transitions
given is fireable from the current state.

\begin{multicols}{2}
\begin{lstlisting}
  boolean-formula =
      ...
    | element deadlock { empty }
    | element is-fireable {
        transition+
      }
    | element integer-le {
        integer-expression,
        integer-expression
      }
    | ...
\end{lstlisting}
\columnbreak
\setlength{\grammarindent}{6em}
\begin{grammar}
<boolean-formula> ::= <atom> | \ldots

<atom> ::= \atomdead
\alt \atomisfire{$t_1, \ldots, t_n$}
\alt <integer-expression> \atomleq <integer-expression>
\end{grammar}
\end{multicols}


\mysubsection{Integer expressions}

%% An integer formula is an integer expression.
%% The tool must return the integer, that is the result of the expression.

%% \begin{multicols}{2}
%% \begin{lstlisting}
%%   integer-formula =
%%       integer-expression

%%   integer-expression = ...
%% \end{lstlisting}
%% \columnbreak
%% \begin{grammar}
%% <integer-formula> ::= <integer-expression>
%% \end{grammar}
%% \end{multicols}

\subsubsection{Arithmetic operators}

These are usual arithmetic operators for integers.

\begin{multicols}{2}
\begin{lstlisting}
  integer-expression =
      ...
    | element integer-constant {
        xsd:integer
      }
    | element integer-sum {
        integer-expression,
        integer-expression+
      }
    | element integer-difference {
        integer-expression,
        integer-expression
      }
    | ...
\end{lstlisting}
\columnbreak
\setlength{\grammarindent}{6em}
\begin{grammar}
<integer-expression> ::=
     \textit{Integer constant}
\alt <integer-expression> \exprplus <integer-expression>
\alt <integer-expression> \exprminus <integer-expression>
\alt \ldots
\end{grammar}
\end{multicols}

\subsubsection{Token-counting integer operator}

\begin{multicols}{2}
\begin{lstlisting}
  integer-expression =
      ...
    | element tokens-count {
        place+
      }
    | ...
\end{lstlisting}
\columnbreak
\begin{grammar}
<integer-expression> ::=
     \atomtokenscnt{$p_1, \ldots, p_n$}
\alt \ldots
\end{grammar}
\end{multicols}

Recall that
\begin{itemize}
  \item \lstinline!tokens-count! returns the exact number of tokens in a set of places.
\end{itemize}

\mysubsection{Bound formulæ}

\begin{multicols}{2}
\begin{lstlisting}
  bound-formula =
      element place-bound {
        place+
      }
\end{lstlisting}
\columnbreak
\begin{grammar}
<bound-formula> ::=
     \atomplacebnd{$p_1, \ldots, p_n$}
\end{grammar}
\end{multicols}

Recall that
\begin{itemize}
\item for a set of places, \lstinline!place-bound! returns some
%interval
%  $[a,b] \subseteq \N$ over the natural numbers,
%  where $a \in \N$ and $b \in \N \cup \set \infty$, if
%  the interval satisfies the following.
%  Let $c$ be the maximum number of tokens that any reachable marking can put in all places of the set.
%  Then~$c$ must be contained in the closed interval $[a,b]$.
natural number $c$ so that $c$ is the maximum number of tokens that any reachable marking can put in all places of the set.
\end{itemize}

\mysubsection{Places and transitions}

Places and transitions are uniquely identified. The identifiers occurring at
the formula XML files are the \texttt{id}'s present in the PNML model.

\begin{lstlisting}
  place =
      element place {
        xsd:IDREF
      }

  transition =
      element transition {
        xsd:IDREF
      }
\end{lstlisting}

\iffalse
\mysection{Parsing the XML formula files}
\label{s:parsing}

The RelaxNG grammar can be downloaded from~\url{http://mcc.lip6.fr/properties/mcc-properties.rnc} using:
\begin{lstlisting}[language=sh]
  wget http://mcc.lip6.fr/properties/mcc-properties.rnc
\end{lstlisting}

The XML schema file can be downloaded from~\url{http://mcc.lip6.fr/properties/mcc-properties.xsd} using:
\begin{lstlisting}[language=sh]
  wget http://mcc.lip6.fr/properties/mcc-properties.xsd
\end{lstlisting}

\mysubsection{How to generate the XML schema from the RelaxNG grammar?}
The Trang tool is able to transform the RelaxNG grammar into an XML Schema.
Visit \url{http://www.thaiopensource.com/relaxng/trang.html} to install this tool.

\begin{lstlisting}[language=sh]
  trang -I rnc -O xsd mcc-properties.rnc mcc-properties.xsd
\end{lstlisting}

\mysubsection{How to generate C++ classes from the XML Schema?}
Generation of {\CC} classes requires Code Synthesis' xsd tool
(\url{http://www.codesynthesis.com/products/xsd/}).
This tool converts the XML Schema of the property language to a set of \CC{} classes,
an XML validating parser, and an XML output.
This tool is free software, and is available for numerous platforms.
It is available at \url{http://www.codesynthesis.com/products/xsd/download.xhtml}.
Parsing and validating XML also requires to install Xerces-\CC{},
available at \url{http://xerces.apache.org/xerces-c/}.

After installing the \lstinline!xsd! tool,
you have to fix the file \lstinline!xsd/cxx/zc-istream.txx!:
\begin{lstlisting}[language=C++]
35c35
<       setg (b, b, e);
---
>       this->setg (b, b, e);
\end{lstlisting}

The conversion from the XML Schema to \CC{} classes is then performed using the following command:
\begin{lstlisting}[language=sh]
  mkdir -o src/cxx/
  xsd cxx-tree \
      --generate-serialization \
      --generate-doxygen \
      --generate-ostream \
      --generate-comparison \
      --generate-detach \
      --generate-default-ctor \
      --generate-polymorphic --polymorphic-type-all \
      --namespace-map http://mcc.lip6.fr=mcc \
      --output-dir src/cxx/ \
      --root-element property-set \
      mcc-properties.xsd
\end{lstlisting}

\mysubsection{How to generate Java classes from the XML Schema?}
Conversion from the XML Schema to Java classes requires the Java Architecture for XML Binding
(JAXB -- \url{http://jaxb.java.net/}).
It is included in recent Java distributions.

To generate the classes, use the following command:
\begin{lstlisting}[language=sh]
  mkdir -o src/java/
  xjc -d src/java/ -p mcc mcc-properties.xsd
\end{lstlisting}
It generates a set of Java files in the \lstinline!java! directory.

\mysubsection{How to generate Python classes from the XML Schema?}
The python script generateDS (\url{http://www.rexx.com/~dkuhlman/generateDS.html}) generates Python code from the XML Schema.

\begin{lstlisting}[language=sh]
  mkdir -p src/python/
  python generateDS.py -m -f --silence -o src/python/mcc-properties.py mcc-properties.xsd
\end{lstlisting}

\mysubsection{How to generate C\# classes from the XML Schema?}
There seems to be also tools for C\# developers: \\
\url{http://stackoverflow.com/questions/386155/comparison-of-xsd-codegenerators-c}.
We did not test them, but are interested by feedback if you use one.
\fi

%% \mysection{Rules Concerning the UpperBounds category}


%% Let $\N \eqdef \set{0, 1, 2, \ldots}$ be the set of natural numbers.
%% %and let $N_n \eqdef \set{0, 1, \ldots, n-1, n}$ be the set including the first
%% %$n+1$ integers.

%% Assume that for a given verification task in the category
%% UpperBounds we have $n \in \N$ answers from $n$ different tools.
%% Let $p_1, \ldots, p_l$ be the concerned net places and let $B$ be the upper
%% bound in the number of tokens present in the $l$ places.
%% Let
%% \[
%% [a_1, b_1], \ldots, [a_n, b_n]
%% \]
%% be the $n$ answers obtained for this verification task.

%% Define function $c \colon \N \to \set{0, \ldots, n}$ as follows.
%% For any $x \in N$, we set $c(x)$ to be the number of tools that returned an
%% interval containing $x$. Formally:
%% \[
%% c(x) \eqdef |\set{i \in \set{1, \ldots, n} \colon x \in [a_i,b_i]}|
%% \]
%% Observe that the intersection of all intervals is non-empty iff
%% there exists some $x \in \N$ such that $c(x) = n$.
%% Let $y \in \N$ be the maximum of the function $c$, that is,
%% $y \eqdef \max\set{c(x) \in \N \colon x \in \N}$.
%% Define now the \emph{master set} $S \subseteq \N$ as follows:
%% \[
%% S \eqdef \set{x \in \N \colon c(x) = y}.
%% \]
%% That is, $S$ is the set of naturals on which the maximum number of tools agree
%% to locate the actual upper bound~$B$.
%% The master set is the best guess that we have to evaluate the $n$ tools and
%% distribute scores. Observe that $S$ can be finite or infinite.

%% Every tool which participated in this verification task is evaluated on three
%% different aspects:
%% \begin{enumerate}
%% \item
%%   How \emph{correct} its answer was.
%% \item
%%   How \emph{precise} its answer was.
%% \item
%%   How \emph{fast} the tool produced an answer.
%% \end{enumerate}
%% The final score is a weighted sum of three sub-scores evaluating these aspects.
%% FIXME.
%% Let $i \in \set{1, \ldots, n}$ be one of the tools.

%% \textbf{Correction of the result}. The sub-score for $i$ is computed as:
%% \begin{itemize}
%% \item
%%   If $S$ is finite, then the sub-score is the
%%   \textit{percentage of inclusion} of $S$ in $[a_i, b_i]$
%%   i.e.\ how much of $S$ is included in $[a_i,b_i]$.
%% \item
%%   If $S$ is infinite, then the sub-score is 0\% or 100\%, depending on whether
%%   the $S$ is entirely included in $[a_i,b_i]$ or not.
%% \end{itemize}

%% \textbf{Precision of the result}. The sub-score for $i$ is computed as:
%% \begin{itemize}
%% \item
%%   If $[a_i,b_i]$ is finite, then the sub-score is the
%%   \textit{percentage of coverage} of $[a_i, b_i]$ in $S$,
%%   i.e.\ how much of $[a_i,b_i]$ is included in~$S$.
%% \item
%%   If $[a_i,b_i]$ is infinite, then the sub-score is 0\% or 100\%, depending on
%%   whether the $[a_i,b_i]$ is entirely included in $S$ or not.
%% \end{itemize}

%% \textbf{Time to compute the result}. FIXME.

\end{document}
